\section{Całka Denjoya w wersji wąskiej}
Absolutna ciągłość jest zbyt silna, aby obejmować funkcje pierwotne wszystkich pochodnych. Wprowadzamy warunek kontrolujący oscylacje na zbiorach, a następnie jego przeliczalną wersję.
\begin{definicja}
Niech $F:[a,b]\to\R$ będzie ciągła i $E\subset[a,b]$. Piszemy $F\in\AC_*(E)$, jeżeli dla każdego $\varepsilon>0$ istnieje $\eta>0$ takie, że dla każdej skończonej rodziny przedziałów $[u_j,v_j]$ o parami rozłącznych wnętrzach i końcach $u_j,v_j\in E$ spełniającej $\sum_j(v_j-u_j)<\eta$ mamy
\[
\sum_j\osc_{[u_j,v_j]}F<\varepsilon.
\]
Piszemy $F\in\AC G_*[a,b]$, jeżeli $[a,b]=\bigcup_{n\ge1}E_n$ dla domkniętych zbiorów $E_n$ takich, że $F\in\AC_*(E_n)$. Na całym przedziale $\AC_*([a,b])=\AC[a,b]$, ale na podzbiorze $E$ warunek oscylacyjny jest silniejszy od kontroli samych przyrostów między punktami $E$.
\end{definicja}
\begin{definicja}[Wąska całka Denjoya]
Funkcja $f:[a,b]\to\R$ o skończonych wartościach jest całkowalna w wąskim sensie Denjoya, jeśli istnieje ciągła $F\in\AC G_*[a,b]$ taka, że $F'=f$ prawie wszędzie. Wówczas definiujemy $(\mathrm D_*)\int_a^b f=F(b)-F(a)$. Jednoznaczność uzasadnimy poniżej.
\end{definicja}
\begin{twierdzenie}[Każda różniczkowalna funkcja należy do $\AC G_*$]
Jeśli $F:[a,b]\to\R$ jest ciągła i ma skończoną pochodną w każdym punkcie wnętrza oraz odpowiednie pochodne jednostronne na końcach, to $F\in\AC G_*[a,b]$.
\end{twierdzenie}
\begin{proof}
Dla $n\in\N$ niech $E_n$ będzie zbiorem tych $x$, dla których $|F(y)-F(x)|\le n|y-x|$ dla wszystkich $y\in[a,b]$ z $|y-x|\le1/n$. Różniczkowalność daje $\bigcup_nE_n=[a,b]$. Każdy $E_n$ jest domknięty: przy $x_k\in E_n$ i $x_k\to x$ przechodzimy do granicy w nierówności, wykorzystując ciągłość $F$. Jeśli $u,v\in E_n$ i $v-u<1/n$, to dla każdego $y\in[u,v]$ mamy $|F(y)-F(u)|\le n(v-u)$. Stąd $\osc_{[u,v]}F\le2n(v-u)$. Dla rodziny przedziałów z końcami w $E_n$ i $\sum_j(v_j-u_j)<\min\{1/n,\varepsilon/(2n)\}$ suma oscylacji jest mniejsza od $\varepsilon$.
\end{proof}
\subsection*{Kontrola oscylacji na zbiorach zerowych}
\begin{lemat}[Zaniedbywalne wahanie na zbiorze zerowym]
\label{it:lem:acstar-null-variation}
Niech $F$ będzie ciągła na $[a,b]$, $E\subset[a,b]$ domknięty i $F\in\AC_*(E)$. Jeżeli $N\subset E$ ma miarę Lebesgue'a zero, to dla każdego $\varepsilon>0$ istnieje cechowanie na $N$, dla którego dowolny skończony drobny podział częściowy $\{(t_j,[u_j,v_j])\}$ ze znacznikami w $N$ spełnia
\[
 \sum_j|F(v_j)-F(u_j)|<\varepsilon.
\]
\end{lemat}
\begin{proof}
Dla pustego $E$ teza jest bezprzedmiotowa. Składowe przedziałowe dopełnienia $E$ w $[a,b]$ są co najwyżej przeliczalne (każda niepusta składowa zawiera liczbę wymierną). Niech $J_k$ oznaczają ich domknięcia. Co najwyżej dwie składowe brzegowe nie mają obu końców w $E$. Suma długości wszystkich $J_k$ nie przekracza $b-a$, ponieważ ich wnętrza są rozłączne.

Twierdzimy, że $\sum_k\osc_{J_k}F<\infty$. Wybieramy $\eta_0>0$ z definicji $\AC_*(E)$ dla błędu $1$. Pomijając skończenie wiele składowych, w tym składowe brzegowe, możemy zapewnić, że suma długości pozostałych jest mniejsza niż $\eta_0$. Każda skończona podrodzina pozostałych ma końce w $E$ i spełnia warunek definicji, więc suma jej oscylacji jest mniejsza od $1$. Supremum po podrodzinach daje zbieżność sumy; skończenie wiele pominiętych oscylacji jest skończonych dzięki ciągłości $F$.

Wybieramy skończoną rodzinę $\mathcal J_0$ składowych (w tym brzegowe) tak, by $\sum_{J_k\notin\mathcal J_0}\osc_{J_k}F<\varepsilon/8$. Niech $m$ oznacza jej liczność. Z definicji $\AC_*(E)$ wybieramy $\eta>0$ dla błędu $\varepsilon/4$. Ponieważ $N$ ma miarę zero, istnieje otwarty zbiór $O\supset N$ o mierze mniejszej od $\eta$. Dla $t\in N$ dobieramy $\delta(t)>0$ tak, by jego otoczenie o promieniu $\delta(t)$ leżało w $O$. Dla każdej składowej $J\in\mathcal J_0$, która nie zawiera $t$, dodatkowo wymagamy $\delta(t)<\operatorname{dist}(t,J)$. Jeżeli $t\in J$, jest jego końcem, gdyż $t\in E$. Z ciągłości $F$ zmniejszamy $\delta(t)$ tak, by $|F(y)-F(t)|<\varepsilon/(8(m+1))$ dla wszystkich $y\in[a,b]$ z $|y-t|<\delta(t)$. Warunki są skończone i wszystkie dopuszczają dodatni promień.

Dla każdej komórki $I_j=[u_j,v_j]$ z rozpatrywanego podziału częściowego zbiór $E\cap I_j$ jest niepusty (zawiera $t_j$) i zwarty. Niech $p_j=\min(E\cap I_j)$, $q_j=\max(E\cap I_j)$. Rozkładamy przyrost:
\[
 |F(v_j)-F(u_j)|
 \le |F(q_j)-F(p_j)|+|F(p_j)-F(u_j)|+|F(v_j)-F(q_j)|.
\]
Przedziały $[p_j,q_j]$ mają rozłączne wnętrza, końce w $E$ i sumę długości nie większą od $\lambda(O)<\eta$. Suma pierwszych składników jest więc mniejsza od $\varepsilon/4$. Jeśli któryś z pozostałych odcinków ma dodatnią długość, leży w domknięciu jednej składowej dopełnienia $E$ i przylega do jej końca w $E$. Dla ustalonej składowej mogą wystąpić najwyżej dwa takie odcinki o dodatniej długości, po jednym przy każdym końcu: dwa przy tym samym końcu miałyby wspólne wnętrze. Składowe brzegowe mają co najwyżej jeden taki odcinek.

Dla składowych spoza $\mathcal J_0$ suma tych składników nie przekracza $2\sum_{J_k\notin\mathcal J_0}\osc_{J_k}F<\varepsilon/4$. Jeżeli komórka wchodzi w składową $J\in\mathcal J_0$, to jej znacznik musi należeć do $J$ (w przeciwnym przypadku warunek odległości wyklucza przecięcie). Znacznik jest więc końcem $J$ w $E$ i odcinek skrajny przy tym końcu ma przyrost mniejszy od $\varepsilon/(8(m+1))$. Łącznie składniki z $\mathcal J_0$ wnoszą mniej niż $\varepsilon/4$. Dodanie trzech oszacowań kończy dowód.
\end{proof}

\begin{wniosek}
\label{it:cor:acgstar-null-variation}
Jeśli $F\in\AC G_*[a,b]$ i $N\subset[a,b]$ ma miarę zero, to wniosek lematu \ref{it:lem:acstar-null-variation} zachodzi dla całego $N$.
\end{wniosek}
\begin{proof}
Niech $[a,b]=\bigcup_{n\ge1}E_n$ będzie domkniętym pokryciem z definicji. Dla $N\cap E_n$ stosujemy lemat z dokładnością $\varepsilon2^{-n-1}$, otrzymując cechowanie $\delta_n$. Każdemu $t\in N$ przypisujemy najmniejsze $n(t)$ z $t\in E_{n(t)}$ i przyjmujemy $\delta(t)=\delta_{n(t)}(t)$. W dowolnym skończonym podziale częściowym grupujemy komórki według $n(t_j)$. Każda grupa spełnia hipotezy odpowiedniego lematu, więc suma modułów przyrostów jest mniejsza od $\sum_{n\ge1}\varepsilon2^{-n-1}=\varepsilon/2$ (dla pustej rodziny suma jest zerowa). W szczególności jest mniejsza od $\varepsilon$.
\end{proof}

\begin{twierdzenie}[Całkowalność pochodnej prawie wszędzie funkcji $\AC G_*$]
\label{it:thm:acgstar-derivative-integrable}
Niech $F\in\AC G_*[a,b]$ oraz niech $f:[a,b]\to\R$ ma skończone wartości i spełnia $F'=f$ prawie wszędzie. Wtedy $f\in\HK[a,b]$ i dla każdego $x\in[a,b]$
\[
 F(x)-F(a)=(\HK)\int_a^x f(t)\,\mathrm dt.
\]
\end{twierdzenie}
\begin{proof}
Zakładamy $a<b$; przypadek przedziału jednopunktowego daje całkę zero. Niech $N$ będzie zbiorem miary zero, poza którym istnieje skończona pochodna równa $f$; dołączamy do $N$ oba końce przedziału. Dla $\varepsilon>0$ wybieramy na $N$ cechowanie z wniosku \ref{it:cor:acgstar-null-variation} dla błędu $\varepsilon/3$. Funkcja $h=f\chi_N$ jest równa zeru poza $N$ i ma skończone wartości. Lemat o zmianach na zbiorze miary zero z rozdziału o całce KH, zastosowany do $|h|$, daje cechowanie dla błędu $\varepsilon/3$. W punktach $N$ bierzemy minimum obu cechowań.

Dla $t\notin N$ z definicji pochodnej dobieramy $\delta(t)>0$, nie większe od cechowania dla $|h|$, aby
\[
 |F(y)-F(t)-f(t)(y-t)|\le
 \frac{\varepsilon}{3(b-a)}|y-t|
 \qquad (|y-t|<\delta(t),\ y\in[a,b]).
\]
Niech $\mathcal P$ będzie dowolnym pełnym podziałem drobnym. Dla komórek ze znacznikami poza $N$ odejmujemy oszacowania dla $y=u_j$ i $y=v_j$; nierówność trójkąta daje błąd nie większy od $\varepsilon(v_j-u_j)/(3(b-a))$. Suma tych błędów nie przekracza $\varepsilon/3$. Dla komórek ze znacznikami w $N$ suma modułów przyrostów $F$ jest mniejsza od $\varepsilon/3$, a suma $|f(t_j)|(v_j-u_j)$ jest mniejsza od $\varepsilon/3$: pełna suma dla $|h|$ jest właśnie tą drugą sumą i ma błąd mniejszy od $\varepsilon/3$ względem całki zero. W konsekwencji
\[
 |S(f,\mathcal P)-F(b)+F(a)|<\varepsilon,
\]
ponieważ suma wszystkich przyrostów $F$ jest teleskopowa. Definicja daje całkowalność i wartość na $[a,b]$. Ten sam argument stosujemy na każdym $[a,x]$: przecięcia zbiorów $E_n$ z podprzedziałem są domkniętym pokryciem, warunek $\AC_*$ pozostaje prawdziwy, a zbiór wyjątkowy pozostaje zerowy po dołączeniu końców. To daje wzór dla wszystkich $x$.
\end{proof}

\subsection*{Różniczkowalność prawie wszędzie}
\begin{lemat}[Przedłużenie afiniczne przez luki]
\label{it:lem:acstar-affine-extension}
Niech $E\subset[a,b]$ będzie niepustym zbiorem domkniętym, a $F:[a,b]\to\R$ funkcją ciągłą należącą do $\AC_*(E)$. Istnieje funkcja $G\in\AC[a,b]$, równa $F$ na $E$, afiniczna na każdej wewnętrznej składowej dopełnienia $E$ i stała na każdej składowej brzegowej dopełnienia.
\end{lemat}
\begin{proof}
Niech $(l_k,r_k)$ będą wewnętrznymi składowymi dopełnienia $E$; obydwa ich końce należą do $E$. Na $E$ przyjmujemy $G=F$, na $[l_k,r_k]$ kładziemy
\[
 G(x)=F(l_k)+\frac{x-l_k}{r_k-l_k}(F(r_k)-F(l_k)).
\]
Na $[a,\min E]$ i $[\max E,b]$ przyjmujemy wartości stałe, odpowiednio $F(\min E)$ i $F(\max E)$. Definicje są zgodne na końcach. Funkcja $G$ jest ciągła: we wnętrzu luk jest afiniczna lub stała, a dla $x\in E$ rozważamy $y\to x$. Jeżeli $y\in E$, korzystamy z ciągłości $F$. Jeżeli $y$ należy do wewnętrznej luki długości mniejszej od $\eta$, jej oba końce są oddalone od $x$ o mniej niż $|x-y|+\eta$, a $G(y)$ jest kombinacją wypukłą ich wartości. Jednostajna ciągłość $F$ pozwala uczynić różnicę $|G(y)-F(x)|$ dowolnie małą przez wybór $\eta$ i $|x-y|$. Luk długości co najmniej $\eta$ jest skończenie wiele; przy $y\to x$ możemy spośród nich rozpatrywać tylko luki, których domknięcia zawierają $x$, i na nich ciągłość w końcu wynika bezpośrednio ze wzoru afinicznego. To samo dotyczy co najwyżej dwóch luk brzegowych.

Z dowodu lematu \ref{it:lem:acstar-null-variation} wiemy, że $\sum_k\osc_{[l_k,r_k]}F<\infty$. W szczególności funkcja
\[
 w(t)=\sum_k\frac{|F(r_k)-F(l_k)|}{r_k-l_k}
             \chi_{(l_k,r_k)}(t)
\]
jest nieujemna, mierzalna i całkowalna w sensie Lebesgue'a: przez zbieżność monotoniczną jej całka jest sumą $|F(r_k)-F(l_k)|$, ograniczoną przez sumę oscylacji. Jeśli $[u,v]$ leży w domknięciu jednej luki, to $|G(v)-G(u)|=\int_u^v w\,\mathrm d\lambda$; w luce brzegowej obie strony są zerowe.

Dla dowolnego $[u,v]$ przecinającego $E$ oznaczamy $p=\min(E\cap[u,v])$, $q=\max(E\cap[u,v])$. Przedziały skrajne $[u,p]$ i $[q,v]$ leżą w lukach lub są jednopunktowe. Zatem
\[
 |G(v)-G(u)|\le |F(q)-F(p)|
       +\int_{[u,p]\cup[q,v]} w\,\mathrm d\lambda.
\]
Rozważmy teraz skończoną rodzinę przedziałów o rozłącznych wnętrzach. Dla tych przecinających $E$ środkowe przedziały $[p,q]$ także mają rozłączne wnętrza i końce w $E$. Przy dostatecznie małej łącznej długości całej rodziny suma $|F(q)-F(p)|$ jest mniejsza od $\varepsilon/2$ przez $\AC_*(E)$. Wszystkie odcinki skrajne i przedziały leżące w lukach mają rozłączne wnętrza, a ich suma długości nie przekracza sumy długości całej rodziny. Absolutna ciągłość całki Lebesgue'a funkcji $w\in L^1$ (rozdział o całce Lebesgue'a) daje dla dostatecznie małej tej sumy łączny wkład całek mniejszy od $\varepsilon/2$. Otrzymujemy warunek absolutnej ciągłości $G$.
\end{proof}

\begin{twierdzenie}[Różniczkowalność prawie wszędzie na zbiorach $\AC_*$]
\label{it:thm:acstar-ae-differentiable}
Jeśli $F$ jest ciągła na $[a,b]$ i $F\in\AC_*(E)$ dla domkniętego $E\subset[a,b]$, to $F$ ma skończoną pochodną prawie wszędzie na $E\cap(a,b)$. W szczególności każda funkcja $F\in\AC G_*[a,b]$ jest różniczkowalna prawie wszędzie.
\end{twierdzenie}
\begin{proof}
Dla pustego $E$ nie ma czego dowodzić. Niech $G$ będzie przedłużeniem z lematu \ref{it:lem:acstar-affine-extension} i $H=F-G$. Na $E$ mamy $H=0$. Ponieważ $G$ jest absolutnie ciągła, jest różniczkowalna prawie wszędzie z podstawowego twierdzenia rachunku całkowego dla całki Lebesgue'a (skrypt \emph{Teoria miary}, rozdział 25, twierdzenie o etykiecie \texttt{thm:FTC-lebesgue-25}). Wystarczy wykazać, że $H'=0$ prawie wszędzie na $E$.

Dla wewnętrznej luki $(l_k,r_k)$ połóżmy $d_k=\osc_{[l_k,r_k]}F$. Z konstrukcji $G$ obie wartości $F(y),G(y)$ leżą pomiędzy minimum i maksimum $F$ na $[l_k,r_k]$. Wobec tego $|H(y)|\le d_k$ dla $y\in[l_k,r_k]$. Ponadto $\sum_k d_k<\infty$ z dowodu lematu \ref{it:lem:acstar-null-variation}.

Dla dodatniej liczby całkowitej $m$ rozważmy zbiory
\[
 U_{k,m}=(l_k-md_k,l_k+md_k)\cup(r_k-md_k,r_k+md_k).
\]
Ich miary nie przekraczają $4md_k$. Jeżeli luk jest nieskończenie wiele, zbiór
\[
 B_m=\bigcap_{n\ge1}\bigcup_{k\ge n}U_{k,m}
\]
ma miarę zero, gdyż dla każdego $n$ jego miara nie przekracza $4m\sum_{k\ge n}d_k$, a prawa strona dąży do zera. Przy skończonej liczbie luk przyjmujemy $B_m=\varnothing$. Wykluczamy z $E$ sumę zbiorów $B_m$, wszystkie końce luk (także brzegowych), oraz $a,b$. Jest to zbiór miary zero.

Ustalmy pozostały punkt $x\in E$ i $m\ge1$. Dla wszystkich dostatecznie dużych $k$ punkt $x$ nie należy do $U_{k,m}$. Jeśli $y$ leży w takiej luce, to jeden z jej końców leży pomiędzy $x$ i $y$, ponieważ $x\in E$. Oznaczając go przez $z$, otrzymujemy $|x-y|\ge |x-z|\ge md_k$. Dla $d_k>0$ mamy więc $|H(y)|/|y-x|\le1/m$; dla $d_k=0$ iloraz jest zerowy. Skończenie wiele pozostałych luk, w tym brzegowe, ma domknięcia niezawierające $x$, zatem można wybrać otoczenie $x$, które nie przecina żadnej z nich. Dla $y\in E$ i $y\ne x$ iloraz jest zerowy. W konsekwencji
\[
 \limsup_{\substack{y\to x\\y\ne x}}
 \frac{|H(y)-H(x)|}{|y-x|}\le\frac1m.
\]
Dowolność $m$ daje $H'(x)=0$. Poza zbiorem zerowym, na którym $G'$ nie istnieje, mamy więc $F'=G'$. Dla funkcji $\AC G_*$ stosujemy wynik do każdego domkniętego zbioru $E_n$ z definicji. Przeliczalna suma zbiorów wyjątkowych jest zbiorem miary zero, a zbiory $E_n$ pokrywają przedział.
\end{proof}

\begin{twierdzenie}[Charakteryzacja funkcji pierwotnych HK; wynik zewnętrzny]
Dla ciągłej $F:[a,b]\to\R$ równoważne są: (i) $F\in\AC G_*[a,b]$; (ii) istnieje $h\in\HK[a,b]$ takie, że $F(x)=F(a)+\int_a^x h$ dla każdego $x$. Wówczas $F'=h$ prawie wszędzie.
\end{twierdzenie}
Twierdzenie w pełnej postaci pozostaje na tym etapie wynikiem zewnętrznym. Twierdzenie \ref{it:thm:acgstar-derivative-integrable} dowodzi już kierunku od funkcji klasy $\AC G_*$ ze znaną pochodną prawie wszędzie do pierwotnej KH. Istnienie pochodnej prawie wszędzie każdej funkcji tej klasy wynika teraz z twierdzenia \ref{it:thm:acstar-ae-differentiable}. Pozostaje udowodnić kierunek od pierwotnej KH do $\AC G_*$; sama ciągłość i pochodna prawie wszędzie nie wystarczają.
\begin{wniosek}[Denjoy = HK]
Wąska całka Denjoya i całka HK mają te same funkcje całkowalne oraz te same wartości. Jeśli $F\in\AC G_*[a,b]$ i $F'=0$ prawie wszędzie, to $F$ jest stała.
\end{wniosek}
\begin{proof}
Dla funkcji HK-całkowalnej jej pierwotna należy do $\AC G_*$ na mocy poprzedniego twierdzenia i ma pochodną równą funkcji podcałkowej prawie wszędzie. Odwrotnie, jeśli $F\in\AC G_*$ oraz $F'=f$ prawie wszędzie, to $F-F(a)=\int_a^x h$ dla pewnego $h\in\HK$, z $h=f$ prawie wszędzie. Lemat o zmianie na zbiorze zerowym pozwala zastąpić $h$ przez $f$. Dla $f=0$ pierwotna jest stała, co zapewnia jednoznaczność. Kombinacje liniowe takich funkcji pierwotnych pozostają w $\AC G_*$.
\end{proof}
\begin{przyklad}[Całkowalność z definicji Denjoya]
Dla $F(x)=x^2\cos(\pi/x^2)$, $x>0$, $F(0)=0$, bierzemy domknięte pokrycie $\{0\},[1/2,1],[1/3,1],\dots$ przedziału $[0,1]$. Na każdym $[1/n,1]$ funkcja $F$ jest klasy $C^1$, więc absolutnie ciągła, a na zbiorze jednopunktowym warunek $\AC_*$ jest automatyczny. Zatem $F\in\AC G_*$ i $F'$ jest całkowalna w sensie Denjoya z wartością $F(1)-F(0)=-1$. Pochodna nie jest jednak absolutnie całkowalna, co wynika z rozbieżności całki modułu jej oscylacyjnego składnika.
\end{przyklad}
\begin{uwaga}[Wąska i szeroka całka Denjoya]
Wąska całka Denjoya ma tę samą klasę funkcji co HK. Szeroką całkę Denjoya opiszemy przez klasę $\AC G$ i pochodną aproksymatywną; jest to istotnie szersza teoria. Historyczna totalizacja Denjoya nie musi odpowiadać w każdym szczególe poniższej definicji opisowej.
\end{uwaga}

% ============================================================
% UZUPELNIENIE MATEMATYCZNE: ROZDZIAL 8
% ============================================================
\subsection*{Jak w praktyce rozpoznawać pierwotne wąskiej całki Denjoya}
\begin{twierdzenie}[Przeliczalny zbiór punktów nieróżniczkowalności]\label{tw:denjoy-countable}
Niech $F:[a,b]\to\R$ będzie ciągła oraz różniczkowalna ze skończoną pochodną w każdym punkcie $(a,b)\setminus N$, gdzie $N\subset[a,b]$ jest zbiorem co najwyżej przeliczalnym. Wtedy $F\in\AC G_*[a,b]$. Jeśli funkcja $f:[a,b]\to\R$ ma skończone wartości i $f=F'$ na $(a,b)\setminus N$, to $f$ jest całkowalna w sensie HK oraz w wąskim sensie Denjoya, a obie całki są równe $F(b)-F(a)$.
\end{twierdzenie}
\begin{proof}
Dla $n\ge1$ definiujemy
\[
 E_n=\left\{x\in[a,b]: |F(y)-F(x)|\le n|y-x|
 \ \text{dla wszystkich}\ y\in[a,b]\ \text{z}\ |y-x|\le1/n\right\}.
\]
Z ciągłości $F$ wynika, że każdy $E_n$ jest domknięty. Każdy punkt, w którym istnieje skończona pochodna (jednostronna na brzegu), należy do któregoś $E_n$. Dla przedziału $[u,v]$ o końcach w $E_n$ i długości mniejszej niż $1/n$ mamy
$\osc_{[u,v]}F\le2n(v-u)$, jak w dowodzie wcześniejszego twierdzenia. Stąd $F\in\AC_*(E_n)$. Zbiory jednopunktowe $\{x\}$, $x\in N\cup\{a,b\}$, są domknięte, a warunek $\AC_*$ zachodzi na nich trywialnie. Przeliczalna rodzina $E_n$ i tych singletonów pokrywa $[a,b]$, więc $F\in\AC G_*$. Ponieważ $N$ ma miarę Lebesgue'a zero, $F'=f$ prawie wszędzie; stosujemy charakteryzację pierwotnych HK i niezmienniczość całki przy zmianie na zbiorze zerowym.
\end{proof}
\begin{przyklad}[Funkcja o przeliczalnie wielu osobliwościach]
Niech $c_n=2^{-n}$, $r_n=2^{-n-3}$ oraz niech $\psi\in C_c^1((-1,1))$ spełnia $\psi(0)=1$. Definiujemy
$F(0)=0$ oraz $F(x)=\sum_{n\ge1}n^{-1}\psi((x-c_n)/r_n)$ dla $x\in(0,1]$. Nośniki składników są parami rozłączne, a amplitudy $1/n$ maleją do zera, więc $F$ jest ciągła na $[0,1]$ i klasy $C^1$ poza punktem $0$. Jej pochodna jest HK-całkowalna oraz $\int_0^1F'=F(1)-F(0)=0$. Ponieważ $\psi$ jest niezerowa i znika na brzegu nośnika, suma wariacji na rozłącznych nośnikach jest nie mniejsza niż $2\sum_{n\ge1}1/n=\infty$. W istocie $\int_0^1|F'(x)|\,\mathrm d\lambda(x)=\left(\int_{-1}^1|\psi'(t)|\,\mathrm dt\right)\sum_{n\ge1}1/n=\infty$. Stąd $F\notin\BV[0,1]$ i $F\notin\AC[0,1]$. Mimo to całka HK pochodnej istnieje.
\end{przyklad}
\begin{stwierdzenie}[Sklejanie pierwotnych]\label{st:denjoy-sklejanie}
Niech $a<c<b$ i niech $F:[a,b]\to\R$ będzie ciągła. Jeśli $F|_{[a,c]}\in\AC G_*[a,c]$ oraz $F|_{[c,b]}\in\AC G_*[c,b]$, to $F\in\AC G_*[a,b]$. Jeśli dodatkowo $F'=f$ prawie wszędzie na obu podprzedziałach, to $(\mathrm D_*)\int_a^b f=(\mathrm D_*)\int_a^c f+(\mathrm D_*)\int_c^b f$.
\end{stwierdzenie}
\begin{proof}
Domknięte pokrycia obu przedziałów przez zbiory, na których obowiązuje $\AC_*$, są także pokryciami przez zbiory domknięte w $[a,b]$. Łącząc te dwie rodziny, otrzymujemy przeliczalne pokrycie $[a,b]$ bez konieczności kontrolowania przedziałów mających końce po różnych stronach $c$: definicja $\AC G_*$ wymaga kontroli osobno na każdym zbiorze pokrycia. Wzór całkowy jest tożsamością przyrostów $F$.
\end{proof}
\begin{uwaga}[Pochodna prawie wszędzie nie identyfikuje pierwotnej]
Ciągłość i równość $F'=f$ prawie wszędzie nie wystarczają do wzoru Newtona--Leibniza. Przykładem jest funkcja Cantora $C$, dla której $C'=0$ prawie wszędzie, lecz $C(1)-C(0)=1$. Warunek $\AC G_*$ usuwa tę niejednoznaczność. Jest zatem warunkiem na \emph{pierwotną}, a nie tylko na zbiór punktów istnienia jej pochodnej.
\end{uwaga}

